\documentclass[10pt, a4paper]{article}

\usepackage{homework}

\title{\textbf{《数值代数》上机作业11}}
\author{岳瀚阳\quad\quad\quad\quad 3230104443}
\date{\today}

\begin{document}

\maketitle

%----------------------------------------------------------------------
\section{共轭梯度法求解 Hilbert-like 矩阵}
%----------------------------------------------------------------------

取矩阵 $A\in\mathbb{R}^{n\times n}$，其元素为
\[
a_{ij} = \frac{1}{i+j+1}, \quad 1\leq i,j\leq n,
\]
右端向量 $b_i = \dfrac{1}{3}\displaystyle\sum_{j=1}^{n}a_{ij}$。则真解为 $x^* = \bigl(\tfrac{1}{3},\ldots,\tfrac{1}{3}\bigr)^\top$，因为
\[
A x^* = \frac{1}{3} A \mathbf{1} = \frac{1}{3}\left[\sum_{j=1}^n a_{ij}\right]_i = b.
\]

$A$ 为对称正定矩阵，故可应用共轭梯度法（CG）。CG 法迭代格式为
\begin{align*}
\alpha_k &= \frac{r_k^\top r_k}{p_k^\top A p_k}, \quad
x_{k+1} = x_k + \alpha_k p_k, \quad
r_{k+1} = r_k - \alpha_k A p_k,\\
\beta_k &= \frac{r_{k+1}^\top r_{k+1}}{r_k^\top r_k}, \quad
p_{k+1} = r_{k+1} + \beta_k p_k,
\end{align*}
停止准则为 $\|r_k\|_2 < 10^{-10}$，初值 $x_0 = 0$。

\subsection{数值结果}

对 $n\in\{5,10,20,50,100\}$ 的数值结果见表~\ref{tab:hilbert}。

\begin{table}[H]
    \centering
    \caption{CG 法求解 Hilbert-like 矩阵的结果}
    \label{tab:hilbert}
    \begin{tabular}{ccccc}
        \toprule
        $n$ & 迭代次数 & $\|x-x^*\|_\infty / \|x^*\|_\infty$ & $\|Ax-b\|_\infty$ & $\kappa(A)$ \\
        \midrule
        5   & 5  & $8.61\times10^{-4}$ & $5.15\times10^{-11}$ & $4.64\times10^{6}$  \\
        10  & 8  & $1.50\times10^{-4}$ & $2.23\times10^{-11}$ & $2.42\times10^{14}$ \\
        20  & 11 & $2.93\times10^{-4}$ & $2.57\times10^{-12}$ & $3.59\times10^{18}$ \\
        50  & 16 & $2.50\times10^{-4}$ & $1.77\times10^{-12}$ & $8.69\times10^{19}$ \\
        100 & 19 & $3.67\times10^{-4}$ & $4.37\times10^{-12}$ & $3.78\times10^{19}$ \\
        \bottomrule
    \end{tabular}
\end{table}

\subsection{结果分析}

\textbf{（1）CG 对所有 $n$ 均收敛。}
增大矩阵维数并不能使 CG 发散。在精确算术下，$n$ 阶对称正定方程组的 CG 法至多 $n$ 步精确收敛，这是理论保证。从表中可见，迭代次数随 $n$ 增大增长极为缓慢（$n=5$ 时 5 步，$n=100$ 时 19 步），远小于上界 $n$。这源于 Hilbert-like 矩阵的特征值高度聚集于 $0$ 附近，仅有少数较大特征值，使 CG 的实际收敛阶远优于 Chebyshev 界
$\left(\dfrac{\sqrt{\kappa}-1}{\sqrt{\kappa}+1}\right)^k$
的预测。

\textbf{（2）条件数饱和于机器精度上限。}
双精度浮点数的机器精度约为 $\varepsilon_{\rm mach} \approx 10^{-16}$，可准确表示的最大条件数约为 $1/\varepsilon_{\rm mach} \approx 10^{16}$。当 $n\geq 12$ 时，$A$ 的真实条件数已远超此值，\texttt{numpy} 计算所得 $\kappa(A)$ 在 $10^{18}\sim10^{20}$ 间波动，仅反映浮点近似的精度上限，而非真实条件数。

\textbf{（3）残差极小，但相对误差停滞。}
残差 $\|Ax-b\|_\infty$ 始终在 $10^{-11}\sim10^{-12}$ 量级（接近机器精度），而相对误差 $\|x-x^*\|_\infty/\|x^*\|_\infty$ 停滞于约 $10^{-4}$，不随 $n$ 改善。这正是极度病态问题的典型现象：$b$ 由同一浮点矩阵 $A$ 计算得到，系统在浮点意义上自洽，CG 能将残差压至机器精度，但解与真值之差受矩阵表示误差支配，精度上限约为 $\varepsilon_{\rm mach}\cdot\kappa(A)\approx10^{-4}$（取 $\kappa=10^{12}$ 的浮点有效分量估算）。

%----------------------------------------------------------------------
\section{三种方法求解 $5\times5$ 方程组}
%----------------------------------------------------------------------

求解方程组 $Ax=b$，其中

\[
A = \begin{bmatrix}
10 &  1 &  2 &  3 &  4 \\
 1 &  9 & -1 &  2 & -3 \\
 2 & -1 &  7 &  3 & -5 \\
 3 &  2 &  3 & 12 & -1 \\
 4 & -3 & -5 & -1 & 15
\end{bmatrix}, \quad
b = \begin{bmatrix}12 \\ -27 \\ 14 \\ -17 \\ 12\end{bmatrix}.
\]

$A$ 为对称矩阵，验证可知 $A$ 正定（各阶顺序主子式均为正：$10,\;89,\;573,\ldots$），故三种方法理论上均可应用。

\subsection{收敛性分析}

将 $A$ 分裂为 $A=D-L-U$，Jacobi 迭代矩阵 $T_J = D^{-1}(L+U)$，G-S 迭代矩阵 $T_{GS}=(D-L)^{-1}U$。数值计算谱半径为
\[
\rho(T_J) = 0.8268, \qquad \rho(T_{GS}) = 0.6846 \approx \rho(T_J)^2.
\]
两者均小于 1，故 Jacobi 与 G-S 均收敛。注意 $A$ 并非严格对角占优（第 1 行和第 3 行的非对角元绝对值之和分别等于、大于对角元），但对称正定条件足以保证收敛。$\rho(T_{GS})\approx\rho(T_J)^2$ 符合 Young 定理在一致次序矩阵下的经典结论。

迭代终止准则：Jacobi 与 G-S 使用 $\|x^{(k+1)}-x^{(k)}\|_\infty < 5\times10^{-5}$，CG 使用 $\|r_k\|_2<10^{-10}$。

\subsection{数值结果}

真解为 $x^*=(1,-2,3,-2,1)^\top$（由 \texttt{np.linalg.solve} 验证）。

\begin{table}[H]
    \centering
    \caption{三种方法求解 $5\times5$ 方程组的结果}
    \label{tab:iter}
    \begin{tabular}{cccc}
        \toprule
        方法 & 迭代次数 & $\|x-x^*\|_\infty/\|x^*\|_\infty$ & $\|Ax-b\|_\infty$ \\
        \midrule
        Jacobi  & 44 & $7.22\times10^{-5}$ & $4.02\times10^{-4}$ \\
        G-S     & 27 & $2.78\times10^{-5}$ & $1.85\times10^{-4}$ \\
        CG      &  5 & $2.59\times10^{-16}$ & $0$ \\
        \bottomrule
    \end{tabular}
\end{table}

\subsection{结果分析}

\textbf{（1）Jacobi 与 G-S 的迭代步数之比约为 $44:27\approx 1.6$。}
因 $\rho(T_{GS})\approx\rho(T_J)^2$，理论上 G-S 所需步数约为 Jacobi 的一半，实测比值略高于 $1/2$，这与停止准则（相邻步差而非对真解的误差）以及非线性截断有关。

\textbf{（2）CG 在 5 步内达到机器精度。}
$A$ 为 $5\times5$ 对称正定矩阵，CG 理论上至多 $n=5$ 步精确收敛；实测恰好 5 步，相对误差仅 $2.6\times10^{-16}$，残差为 $0$，验证了理论上限的精确性。这一性能远优于 Jacobi（44步）和 G-S（27步）：对于良态对称正定系统，CG 无论从收敛速度还是最终精度均具有明显优势。

\textbf{（3）精度差异的本质原因。}
Jacobi 和 G-S 的停止准则 $\|x^{(k+1)}-x^{(k)}\|_\infty<5\times10^{-5}$ 并不直接控制对真解的误差，其最终精度约为 $10^{-4}\sim10^{-5}$。CG 的精度受限于机器精度而非迭代截断，因此对于维数较小、条件数适中的对称正定系统，CG 是三种方法中最优的选择。

%----------------------------------------------------------------------
\appendix
\section{源代码}
%----------------------------------------------------------------------

\subsection{\texttt{hw11\_test.py}（主程序）}
\lstinputlisting{../tests/hw11_test.py}

\subsection{\texttt{conjugate\_gradient.py}（共轭梯度法）}
\lstinputlisting{../algorithms/conjugate_gradient.py}

\subsection{\texttt{iteration.py}（Jacobi / G-S 迭代法）}
\lstinputlisting{../algorithms/iteration.py}

\end{document}
